\chapter{Background e stato dell'arte}
\label{cap:02-stato-dell-arte}

Questo capitolo fissa il vocabolario e colloca \nf{} rispetto al lavoro
esistente. Non \`e una rassegna esaustiva del serverless --- ne esistono di
migliori~\cite{jonas2019berkeley, baldini2017trends} --- ma una ricostruzione
mirata delle quattro linee di ricerca su cui il progetto interviene: la
struttura delle piattaforme FaaS, il costo di avvio a freddo, le politiche di
elasticit\`a e di ammissione, e l'offloading fra dispositivi periferici e
cloud.

\section{Una definizione operativa di FaaS}

Nella letteratura il termine \emph{serverless} copre sia un modello economico
(pagamento a consumo, granularit\`a fine, nessun costo a riposo) sia un modello
di programmazione. Qui interessa il secondo, e conviene definirlo per ci\`o che
richiede alla piattaforma anzich\'e per ci\`o che promette all'utente. Una
piattaforma FaaS \`e un sistema che soddisfa quattro obblighi:

\begin{enumerate}[leftmargin=*]
  \item \textbf{Registrazione.} Accetta la definizione di un'unit\`a di calcolo
  --- un'immagine, un comando, un insieme di variabili d'ambiente, dei limiti di
  risorsa --- e le associa un nome logico stabile.
  \item \textbf{Attivazione.} Su richiesta, garantisce l'esistenza di almeno
  un'istanza in esecuzione capace di servire quel nome, creandola se assente.
  \item \textbf{Instradamento.} Consegna una richiesta a una delle istanze e ne
  restituisce l'esito al chiamante, rispettando un contratto di
  serializzazione noto.
  \item \textbf{Ritiro.} Riduce il numero di istanze quando la domanda cala,
  eventualmente fino a zero.
\end{enumerate}

Sono obblighi minimi: tutto ci\`o che le piattaforme reali offrono in pi\`u ---
composizione di funzioni, trigger da sorgenti di eventi, versioning, traffic
splitting, isolamento multi-tenant --- \`e sovrastruttura rispetto a questo
nucleo. \nf{} implementa i quattro obblighi e quasi nient'altro; il
\cref{cap:03-requisiti-principi} argomenta perch\'e.

Va osservato che gli obblighi 2 e 4 sono in tensione fra loro, ed \`e da questa
tensione che discende la maggior parte dei problemi interessanti. Ritirare
istanze riduce il costo a riposo ma allunga l'attivazione successiva; mantenerle
accorcia l'attivazione ma paga capacit\`a inutilizzata. Il punto in cui una
piattaforma colloca il proprio compromesso fra i due \`e, di fatto, la sua
identit\`a.

\section{Anatomia delle piattaforme di riferimento}

Le piattaforme FaaS open source per Kubernetes condividono una struttura
ricorrente, che vale la pena esplicitare perch\'e \`e precisamente la struttura
da cui \nf{} si discosta.

\paragraph{Knative Serving~\cite{knative}.} \`E la realizzazione pi\`u
articolata del modello. Una richiesta entra da un ingress di rete, raggiunge un
componente chiamato \emph{activator} quando la revisione ha zero repliche,
attende che l'\emph{autoscaler} (KPA) osservi la domanda e comandi allo
scheduler di Kubernetes la creazione di pod, e viene infine instradata a un
sidecar (\emph{queue-proxy}) che, all'interno del pod applicativo, applica il
limite di concorrenza e riporta le metriche all'autoscaler. Il modello di
concorrenza \`e esplicito e centrale: la grandezza \code{containerConcurrency}
governa contemporaneamente l'ammissione locale e il segnale di scaling. La
qualit\`a del progetto \`e alta e il numero di componenti che una singola
invocazione attraversa \`e, di conseguenza, elevato.

\paragraph{Apache OpenWhisk~\cite{openwhisk}.} Adotta un'architettura a
\emph{controller} e \emph{invoker}: il controller riceve l'attivazione, consulta
un database (CouchDB) per la definizione dell'azione, e la accoda su Kafka verso
un invoker, che mantiene un pool di container caldi e ne riusa uno se
disponibile. La distinzione fra container \emph{cold}, \emph{prewarm} e
\emph{warm} \`e resa esplicita nel progetto, e la letteratura sul cold start ha
spesso usato OpenWhisk come piattaforma di riferimento~\cite{wang2018peeking}.
Il costo \`e la dipendenza da due sistemi di stato esterni sul percorso di ogni
invocazione.

\paragraph{OpenFaaS~\cite{openfaas}.} Semplifica considerevolmente: un gateway
riceve la richiesta e la inoltra a un Service Kubernetes che bilancia sulle
repliche del Deployment della funzione; lo scaling \`e delegato a componenti
esterni (HPA o KEDA). Il percorso sincrono \`e corto, e l'asincrono passa per
una coda (NATS). \`E la piattaforma pi\`u vicina a \nf{} per filosofia, ma
delega all'esterno esattamente le politiche che \nf{} vuole rendere
sostituibili.

\paragraph{Fission~\cite{fission}.} Introduce il concetto di \emph{pool} di pod
generici pre-avviati in cui il codice della funzione viene iniettato al momento
dell'uso, riducendo drasticamente il tempo di attivazione a scapito
dell'isolamento e della generalit\`a del runtime.

\Cref{tab:piattaforme} riassume il confronto. La colonna che conta ai fini di
questo lavoro \`e l'ultima: il numero di processi distinti che una singola
invocazione sincrona attraversa nel caso caldo.

\begin{table}[htbp]
\centering
\caption{Confronto strutturale fra piattaforme FaaS per Kubernetes.
Il conteggio dei salti \`e indicativo e riferito al percorso sincrono con
repliche gi\`a calde, escluso l'ingress di rete.}
\label{tab:piattaforme}
\small
\begin{tabularx}{\textwidth}{@{}lXccc@{}}
\toprule
\textbf{Piattaforma} & \textbf{Componenti sul percorso} & \textbf{Stato esterno} &
\textbf{Scale-to-zero} & \textbf{Salti} \\
\midrule
Knative     & ingress, activator, queue-proxy, utente & no & s\`i & 3--4 \\
OpenWhisk   & controller, Kafka, invoker, utente      & CouchDB, Kafka & s\`i & 3--4 \\
OpenFaaS    & gateway, Service, utente                & NATS (async) & con add-on & 2 \\
Fission     & router, executor, pod generico, utente  & etcd & s\`i & 2--3 \\
\nf{}       & control plane, utente                   & nessuno & no & 1 \\
\bottomrule
\end{tabularx}
\end{table}

La riga di \nf{} non \`e un vanto prestazionale: un singolo salto \`e
raggiungibile perch\'e la piattaforma rinuncia a tutto ci\`o che gli altri salti
fornivano. \`E per\`o la condizione che rende possibile l'esperimento descritto
nel \cref{cap:25-valutazione-concorrenza}, in cui una variazione di alcuni
millisecondi nel tempo di attesa dev'essere attribuita a una politica di
ammissione e non a una catena di proxy interposti.

\section{Il cold start e le sue mitigazioni}

L'attivazione di una funzione che non ha istanze pronte impone al chiamante un
ritardo composto da pi\`u contributi: la decisione di scheduling, il pull
dell'immagine, la creazione del sandbox, l'inizializzazione del runtime (che per
una JVM o un interprete Python con dipendenze pesanti domina spesso tutti gli
altri) e infine l'esecuzione dell'handler.

La letteratura ha attaccato ciascun contributo separatamente.

\begin{itemize}[leftmargin=*]
  \item \textbf{Sandbox pi\`u leggeri.} Firecracker~\cite{agache2020firecracker}
  dimostra che una microVM con superficie hardware ridotta pu\`o avviarsi in
  poche decine di millisecondi mantenendo isolamento a livello di hypervisor.
  SOCK~\cite{oakes2018sock} riprogetta il container per il caso serverless,
  eliminando dal percorso di avvio le operazioni del runtime container generico
  che il caso d'uso non richiede.

  \item \textbf{Riuso dello stato inizializzato.}
  Catalyzer~\cite{du2020catalyzer} e i lavori su snapshot e
  restore~\cite{ustiugov2021snapshots} sostituiscono l'inizializzazione con il
  ripristino di un'immagine di memoria gi\`a inizializzata, spostando il costo
  dal percorso critico a una fase preparatoria.

  \item \textbf{Politiche di ritenzione.} FaasCache~\cite{fuerst2021faascache}
  riformula il mantenimento dei container caldi come un problema di caching e
  applica politiche di sostituzione greedy-dual; \cite{shahrad2020wild} mostra su
  traccia di produzione che le distribuzioni di inter-arrivo reali sono
  sufficientemente strutturate da rendere efficace la predizione.
\end{itemize}

\nf{} non contribuisce a questa linea. Misura il cold start
(\code{function\_cold\_start\_ms}, \cref{cap:21-osservabilita}) e lo tratta come
una grandezza osservabile del proprio ambiente, ma non introduce meccanismi per
ridurlo. La scelta \`e coerente con il posizionamento del progetto: il cold
start su \nf{} \`e dominato dal tempo di creazione del pod da parte di
Kubernetes e dall'avvio del runtime della funzione, entrambi fuori dal
controllo del control plane.

\section{Elasticit\`a: quante repliche}

Il problema dell'autoscaling di applicazioni elastiche \`e antecedente al
serverless e ne esiste una rassegna
consolidata~\cite{loridobotran2014autoscaling}, che classifica gli approcci in
reattivi a soglia, basati su teoria del controllo, su modelli di code, su
apprendimento per rinforzo e su serie temporali.

Nel contesto Kubernetes le due realizzazioni dominanti sono reattive a soglia.
L'Horizontal Pod Autoscaler~\cite{k8shpa} confronta periodicamente un valore
osservato (tipicamente l'utilizzazione di CPU, oppure una metrica custom) con
un valore obiettivo e calcola il numero di repliche come rapporto fra i due,
applicando bande morte e finestre di stabilizzazione per evitare oscillazioni.
KEDA~\cite{keda} generalizza la sorgente della metrica a un catalogo esteso di
\emph{scaler} orientati agli eventi, e --- soprattutto --- consente lo scaling
a zero, che l'HPA di base non permette.

Knative adotta un autoscaler dedicato la cui metrica primaria non \`e
l'utilizzazione di CPU ma la \emph{concorrenza media per replica}, ossia il
numero medio di richieste contemporaneamente in carico su ciascun pod. \`E una
scelta significativa, perch\'e lega la decisione di scaling alla stessa
grandezza che governa l'ammissione locale.

\nf{} riproduce questa impostazione nel modulo \code{autoscaler}
(\cref{cap:11-autoscaler}), che pu\`o decidere da richieste al secondo,
profondit\`a di coda o richieste in volo. La differenza rilevante \`e che \nf{}
tratta il numero di repliche e il limite di concorrenza per funzione come due
decisioni \emph{separabili e separatamente sostituibili}, mentre nelle
piattaforme citate sono accoppiate nello stesso componente.

\section{Ammissione, backpressure e controllo della concorrenza}

Quante richieste ammettere contemporaneamente \`e un problema pi\`u antico del
serverless e con una letteratura pi\`u matura.

\paragraph{Il risultato strutturale.} La legge di Little~\cite{little1961proof}
lega il numero medio di richieste presenti nel sistema $L$, il tasso di arrivo
$\lambda$ e il tempo medio di permanenza $W$ dalla relazione $L = \lambda W$.
Applicata a un sistema con limite di concorrenza $C$ saturo, fornisce
immediatamente $W = C / \lambda$: fissato il tasso di arrivo, il tempo di
permanenza \`e proporzionale al limite. Questa relazione, banale in s\'e, ha una
conseguenza che il \cref{cap:12-concurrency-control} sviluppa in dettaglio: un
controllore che aumenta $C$ per aumentare il throughput sta necessariamente
allungando la permanenza, e i due obiettivi non possono essere ottimizzati
indipendentemente. La teoria delle code~\cite{kleinrock1975queueing} fornisce
l'apparato per rendere quantitativa l'osservazione.

\paragraph{Architetture a stadi.} SEDA~\cite{welsh2001seda} propone di
strutturare un servizio come una pipeline di stadi separati da code esplicite,
ciascuno con il proprio pool di thread dimensionato dinamicamente. Il contributo
concettuale rilevante qui \`e l'idea che le code \emph{debbano essere esplicite e
osservabili}: un servizio ben condizionato \`e uno in cui il punto in cui le
richieste si accumulano \`e noto e strumentato, anzich\'e disperso nei buffer
del sistema operativo. Il modulo \code{sync-queue} di \nf{}
(\cref{cap:10-sync-queue}) segue direttamente questa impostazione.

\paragraph{Controllo dal ritardo.} La linea di lavoro sul controllo di
congestione fornisce due meccanismi che \nf{} adatta.
TCP Vegas~\cite{brakmo1995vegas} stima il ritardo di coda confrontando il
round-trip time corrente con il minimo osservato --- assunto come stima del
ritardo di propagazione puro --- e regola la finestra di trasmissione in base
alla differenza, anzich\'e attendere la perdita di pacchetti come segnale.
CoDel~\cite{nichols2012codel} osserva che il segnale utile non \`e la lunghezza
della coda ma il \emph{ritardo minimo sperimentato in una finestra temporale},
e scarta pacchetti quando tale minimo resta sopra soglia per un intervallo
prolungato. Entrambi condividono la stessa intuizione: il ritardo di coda \`e
distinguibile dal costo di servizio irriducibile solo confrontando l'osservazione
corrente con un riferimento storico.

La libreria \code{concurrency-limits} di Netflix~\cite{netflixlimits} porta
esplicitamente questi algoritmi nel dominio delle chiamate RPC, implementando
limitatori in stile Vegas e Gradient che aggiustano un limite di concorrenza
osservando la latenza. \nf{} riprende il segnale in stile Vegas nel modo
\code{ADAPTIVE\_PER\_POD} (\cref{cap:12-concurrency-control}), e il
\cref{cap:25-valutazione-concorrenza} riporta il motivo per cui, nel contesto
FaaS, questo segnale si \`e rivelato insufficiente.

\paragraph{Equit\`a nella ripartizione.} Quando pi\`u flussi competono per una
risorsa condivisa e la somma delle richieste eccede la disponibilit\`a, la
nozione di allocazione max-min~\cite{bertsekas1992datanetworks} fornisce un
criterio: si massimizza l'allocazione minima, e chi chiede meno della propria
quota equa la ottiene per intero, mentre l'eccedenza viene ridistribuita fra i
restanti. Nella variante pesata ciascun flusso porta un peso che scala la
propria quota. Il modo \code{BUDGETED} di \nf{} applica questo criterio ai limiti
di concorrenza per funzione, trattando la capacit\`a della piattaforma come la
risorsa condivisa e il fabbisogno stimato di ciascuna funzione come la sua
domanda.

\section{Offloading fra periferia e cloud}

L'idea di spostare il calcolo dal dispositivo o dal nodo periferico verso
un'infrastruttura pi\`u capiente quando la periferia non basta risale a prima
del serverless~\cite{bonomi2012fog, shi2016edge, satya2017edge}. La formulazione
classica \`e un problema di decisione: dato un compito, valutare se eseguirlo
localmente o trasferirlo, sulla base di un modello del costo di trasferimento e
del guadagno di esecuzione.

La declinazione che \nf{} adotta nel modulo \code{offload}
(\cref{cap:13-offload}) \`e volutamente pi\`u modesta e pi\`u operativa. La
decisione non \`e presa da un modello ma da due segnali osservati: una politica
per funzione dichiarata dall'operatore (\emph{eager}), oppure il rifiuto della
coda di ammissione locale per profondit\`a o per attesa stimata eccessiva
(\emph{pressure}). Il trasferimento \`e a singolo salto e senza fallback locale.
Questa combinazione --- decidere dal rifiuto anzich\'e da un modello di costo ---
non risolve il problema di ottimizzazione classico, ma ha la propriet\`a di non
richiedere alcuna calibrazione preventiva, e di degradare in modo prevedibile.

\section{Benchmarking delle piattaforme serverless}

Confrontare piattaforme FaaS \`e notoriamente difficile: il risultato dipende dal
carico, dalla dimensione del payload, dal profilo di arrivo, dall'hardware, e
dallo stato caldo o freddo del sistema al momento della misura.
SeBS~\cite{copik2021sebs} e ServerlessBench~\cite{yu2020serverlessbench}
propongono suite di carichi standardizzati e metriche comparabili, e
\cite{wang2018peeking} mostra quanto informazione si possa estrarre sulle
politiche interne di una piattaforma commerciale per sola osservazione esterna.

Il metodo sperimentale di \nf{} (\cref{cap:22-metodo-sperimentale}) non aspira
alla comparabilit\`a inter-piattaforma. Persegue invece la comparabilit\`a
\emph{intra}-piattaforma nel tempo: fissato un profilo hardware e un carico, i
risultati di ogni release sono confrontati con quelli della release precedente
sullo stesso profilo, e la deviazione \`e trattata come un segnale di
regressione. \`E un obiettivo pi\`u ristretto e pi\`u difendibile.

\section{Posizionamento}

Riassumendo, \nf{} non contende alle piattaforme mature il terreno della
completezza funzionale, e non contribuisce alla riduzione del cold start. Si
colloca invece nello spazio, meno affollato, di una piattaforma FaaS
sufficientemente piccola da essere usata come banco di prova per politiche di
ammissione, scaling e concorrenza, con tre elementi di novit\`a rispetto a
quanto sopra:

\begin{enumerate}[leftmargin=*]
  \item la \textbf{separazione esplicita} fra la decisione sul numero di repliche
  e la decisione sul limite di concorrenza per funzione, mantenute in due moduli
  indipendenti e sostituibili separatamente;
  \item la formulazione di un controllore che decide dal \textbf{tempo di
  sojourn} anzich\'e dal tempo di servizio, con la dimostrazione --- analitica e
  sperimentale --- che la seconda scelta non ammette ottimo interno;
  \item un modo di allocazione \textbf{a budget di piattaforma} che tratta la
  concorrenza come risorsa condivisa soggetta a equit\`a max-min pesata, anzich\'e
  come parametro che ogni funzione negozia per conto proprio.
\end{enumerate}
